\documentclass{article}
\usepackage{amsmath, amssymb, amsthm}
\usepackage{hyperref}
\usepackage{geometry}
\geometry{margin=1in}

\title{Erd\H{o}s Problem \#668}
\date{Last updated: 2025-08-31}
\author{Source: erdosproblems.com}

\begin{document}
\maketitle

\noindent\textbf{Status:} OPEN \quad \textbf{No prize}

\noindent\textbf{Tags:} geometry, distances

\medskip

\noindent\textbf{Problem Statement:}

Is it true that the number of incongruent sets of $n$ points in $\mathbb{R}^2$ which maximise the number of unit distances tends to infinity as $n\to\infty$? Is it always $&#62;1$ for $n&#62;3$?




In fact this is $=1$ also for $n=4$, the unique example given by two equilateral triangles joined by an edge. Computational evidence of Engel, Hammond-Lee, Su, Varga, and Zs\'{a}mboki \cite{EHSVZ25} and Alexeev, Mixon, and Parshall \cite{AMP25} suggests that this count is $=1$ for various other $5\leq n\leq 21$ (although these calculations were checking only up to graph isomorphism, rather than congruency). The actual maximal number of unit distances is the subject of [90] .




References

[AMP25] B. Alexeev, D. Mixon, and H. Parshall, The Erd\H{o}s unit distance problem for small point sets . arXiv:2412.11914 (2025).

[EHSVZ25] P. Engel, O. Hammond-Lee, Y. Su, D. Varga, and P. Zs\'{a}mboki, Diverse beam search to find densest-known planar unit distance graphs . arXiv:2406.15317 (2025).

\noindent\textbf{Related OEIS sequences:} A385657


\bigskip
\noindent\small{Source: \url{https://www.erdosproblems.com/668}}
\end{document}
